We want to generalise this result to any representation-infinite poset. We recall that a given representation-infinite poset $P$ can be reduced to a critical poset $C$ (by Lemma \ref{lemma:infinitereduction}). We just proved that $\incidencealgebra{C}$ is brick-continuous and we constructed a generic brick for it. If we can ``reverse'' the reduction process we may be able to compute a generic brick for all the representation-infinite posets. For this purpose, we will now define the ``extension'' functors of the contraction and of the subtraction and prove that they have the necessary properties.

\subsubsection{Extension functor of contraction}\label{subsection:extfuncontr}

We recall that any contraction of finite posets can be written as a composite of elementary contractions, so we will focus our attention on elementary contractions.

Let $P$, $P'$ be finite posets with $c\colon P \to P'$ an elementary contraction. If $\card{P}=n$ then $\card{P'}=n-1$ (since $c$ is elementary), so, up to re-indexing the elements of the posets, we may suppose that $P=\{x_1,\dots,x_{n-1},x_n\}$ with $x_{n-1}\leq x_n$ neighbours and $P'=\{y_1,\dots,y_{n-1}\}$ and $c$ defined by $x_i \mapsto y_i$ for $i=1,\dots, n-2$ and $x_{n-1},x_n \mapsto y_{n-1}$.
Let $(\HasseQuiver{P}, I_{P}), (\HasseQuiver{P'}, I_{P'})$ be the Hasse quivers with relations of $P'$ and $P'$ respectively.

Consider the natural map of right projective $\incidencealgebra{P}$-modules $\varphi\colon e_{n-1}(\incidencealgebra{P}) \to e_n(\incidencealgebra{P})$.
We have a fully faithful functor $G\defeq g_*\colon \Mod{(\incidencealgebra{P})_\varphi} \hookrightarrow \Mod{\incidencealgebra{P}}$ given by the universal localisation $g \colon \incidencealgebra{P}\to (\incidencealgebra{P})_\varphi$. Moreover we know that (see \cite[Theorem 4.7]{schofield1985representation})
\[
\Mod{(\incidencealgebra{P})_\varphi} \cong \{M\in \Rep{\HasseQuiver{P}, I_{P}} \sth \text{the map } M_{n-1}\to M_n \text{ is an isomorphism}\}
\]
We get an equivalence $F\colon \Mod{\incidencealgebra{P'}} \xrightarrow{\cong} \Mod{(\incidencealgebra{P})_\varphi}$ in the following way: for a representation $M\in \Rep{\HasseQuiver{P'}, I_{P'}}\cong \Mod{\incidencealgebra{P'}}$ we define 
\[
F(M)_i \defeq 
\begin{cases}
	M_i &\text{for } i=1,\dots, n-1 \\
	M_{n-1} &\text{for } i=n
\end{cases}
\]
and we set $M_{n-1} = F(M)_{n-1} \xrightarrow{\text{id}_{M_{n-1}}} F(M)_n = M_{n-1}$. The other $\K$-linear maps in $F(M)$ are defined from those in $M$ in order to satisfy the relations in $I_P$. For example
% https://q.uiver.app/#q=WzAsMTgsWzEsMywiMSJdLFswLDIsIjIiXSxbMSwxLCI0Il0sWzIsMiwiMyJdLFsxLDAsIlAiXSxbNCwzLCIxIl0sWzQsMiwiMiJdLFs0LDEsIjMiXSxbNCwwLCJQJyJdLFs2LDMsIlZfMSJdLFs2LDIsIlZfMiJdLFs2LDEsIlZfMyJdLFs4LDMsIlZfMSJdLFs3LDIsIlZfMiJdLFs4LDEsIlZfMyJdLFs5LDIsIlZfMyJdLFs2LDAsIlxcTW9ke1xcaW5jaWRlbmNlYWxnZWJyYXtcXEt9e1AnfX0iXSxbOCwwLCJcXE1vZHsoXFxpbmNpZGVuY2VhbGdlYnJhe1xcS317UH0pX1xcdmFycGhpfSJdLFswLDFdLFsxLDJdLFswLDNdLFszLDJdLFswLDIsIiIsMSx7InN0eWxlIjp7ImJvZHkiOnsibmFtZSI6ImRhc2hlZCJ9LCJoZWFkIjp7Im5hbWUiOiJub25lIn19fV0sWzUsNl0sWzYsN10sWzQsOCwiYyIsMCx7InN0eWxlIjp7ImhlYWQiOnsibmFtZSI6ImVwaSJ9fX1dLFs5LDEwLCJcXHBzaV97MTJ9IiwyXSxbMTAsMTEsIlxccHNpX3syM30iLDJdLFsxMiwxMywiXFxwc2lfezEyfSJdLFsxMywxNCwiXFxwc2lfezIzfSJdLFsxMiwxNSwiXFxwc2lfezIzfVxccHNpX3sxMn0iLDJdLFsxNSwxNCwiMSIsMl0sWzEyLDE0LCIiLDEseyJzdHlsZSI6eyJib2R5Ijp7Im5hbWUiOiJkYXNoZWQifSwiaGVhZCI6eyJuYW1lIjoibm9uZSJ9fX1dLFsxNiwxNywiRiJdXQ==
\[\begin{tikzcd}[ampersand replacement=\&,cramped,sep=small]
	\& P \&\&\& {P'} \&\& {\Mod{\incidencealgebra{P'}}} \&\& {\Mod{(\incidencealgebra{P})_\varphi}} \& \\
	\& 4 \&\&\& 3 \&\& {V_3} \&\& {V_3} \\
	2 \&\& 3 \&\& 2 \&\& {V_2} \& {V_2} \&\& {V_3} \\
	\& 1 \&\&\& 1 \&\& {V_1} \&\& {V_1}
	\arrow["c", two heads, from=1-2, to=1-5]
	\arrow["F", from=1-7, to=1-9]
	\arrow[from=3-1, to=2-2]
	\arrow[from=3-3, to=2-2]
	\arrow[from=3-5, to=2-5]
	\arrow["{\psi_{23}}"', from=3-7, to=2-7]
	\arrow["{\psi_{23}}", from=3-8, to=2-9]
	\arrow["1"', from=3-10, to=2-9]
	\arrow[dashed, no head, from=4-2, to=2-2]
	\arrow[from=4-2, to=3-1]
	\arrow[from=4-2, to=3-3]
	\arrow[from=4-5, to=3-5]
	\arrow["{\psi_{12}}"', from=4-7, to=3-7]
	\arrow[dashed, no head, from=4-9, to=2-9]
	\arrow["{\psi_{12}}", from=4-9, to=3-8]
	\arrow["{\psi_{23}\psi_{12}}"', from=4-9, to=3-10]
\end{tikzcd}\]


Using the previous description of $\Mod{(\incidencealgebra{P})_\varphi}$, it can be easily seen that $F$ is a well-defined fully faithful dense functor, giving us an equivalence.


\begin{definition}
	We define the \textbf{extension functor} $E_c$ of the elementary contraction $c$ as the fully faithful functor $E_c \defeq G\circ F \colon \Mod{\incidencealgebra{P'}} \xrightarrow{\cong} \Mod{(\incidencealgebra{P})_\varphi} \hookrightarrow \Mod{\incidencealgebra{P}}$.
\end{definition}

By the definition of the functors $F$ and $G$, it is clear that the restriction $E_{c \downarrow} \colon \Modf{\incidencealgebra{P'}} \hookrightarrow \Modf{\incidencealgebra{P}}$ is also a well-defined fully faithful functor.

%%%%%%%%%%%%%%%%%%%%%%%%%%%%%%%%%%%%%%%%%%%%%%%%%%%%%%%%%%%%%%

\subsubsection{Extension functors of subtraction}
We recall that given a subtraction of a poset we may suppose to remove one vertex at the time (c.f. Definition \ref{def:reduction}), so we will focus our attention on subtractions that remove just one element.

Let $P$ be a finite poset, $P'$ a subposet of $P$ with $n=\card{P}=\card{P'}+1$ and denote by $s\colon P' \to P$ the inclusion of $P'$ in $P$ as a subposet. Up to re-indexing the elements of the posets we may suppose that $P=\{1,\dots,n-1, n\}$ and $P'=\{1,\dots,n-1\}$, so the subtraction removed the vertex $n$.
Let $(\HasseQuiver{P}, I_{P}), (\HasseQuiver{P'}, I_{P'})$ be the bound Hasse quivers of $P$ and $P'$ respectively. It is easy to check that $\incidencealgebra{P'} \cong e(\incidencealgebra{P})e$ where $e=\sum_{i=1}^{n-1} e_i$ with $e_i$ the idempotent corresponding to the vertex $i$. Therefore we have the following $\K$-linear covariant functors, with their restrictions to finite-dimensional modules (c.f.~\cite[Ch. I.6]{sim:vol1})
% https://q.uiver.app/#q=WzAsNCxbMCwwLCJcXE1vZHtcXGluY2lkZW5jZWFsZ2VicmF7XFxLfXtQfX0iXSxbMiwwLCJcXE1vZHtcXGluY2lkZW5jZWFsZ2VicmF7XFxLfXtQJ319Il0sWzQsMCwiXFxNb2Rme1xcaW5jaWRlbmNlYWxnZWJyYXtcXEt9e1B9fSJdLFs2LDAsIlxcTW9kZntcXGluY2lkZW5jZWFsZ2VicmF7XFxLfXtQJ319Il0sWzAsMSwiXFx0ZXh0e3Jlc31fZSJdLFsxLDAsIlRfZSIsMix7Im9mZnNldCI6LTMsImN1cnZlIjo0LCJzaG9ydGVuIjp7InNvdXJjZSI6MTAsInRhcmdldCI6MTB9fV0sWzEsMCwiTF9lIiwwLHsib2Zmc2V0IjozLCJjdXJ2ZSI6LTQsInNob3J0ZW4iOnsic291cmNlIjoxMCwidGFyZ2V0IjoxMH19XSxbMiwzLCJcXHRleHR7cmVzfV9lIl0sWzMsMiwiVF9lIiwyLHsib2Zmc2V0IjotMywiY3VydmUiOjQsInNob3J0ZW4iOnsic291cmNlIjoxMCwidGFyZ2V0IjoxMH19XSxbMywyLCJMX2UiLDAseyJvZmZzZXQiOjMsImN1cnZlIjotNCwic2hvcnRlbiI6eyJzb3VyY2UiOjEwLCJ0YXJnZXQiOjEwfX1dXQ==
\[\begin{tikzcd}[ampersand replacement=\&,cramped]
	{\Mod{\incidencealgebra{P}}} \&\& {\Mod{\incidencealgebra{P'}}} \&\& {\Modf{\incidencealgebra{P}}} \&\& {\Modf{\incidencealgebra{P'}}}
	\arrow["{\text{res}_e}", from=1-1, to=1-3]
	\arrow["{T_e}"', shift left=3, curve={height=24pt}, between={0.1}{0.9}, from=1-3, to=1-1]
	\arrow["{L_e}", shift right=3, curve={height=-24pt}, between={0.1}{0.9}, from=1-3, to=1-1]
	\arrow["{\text{res}_e}", from=1-5, to=1-7]
	\arrow["{T_e}"', shift left=3, curve={height=24pt}, between={0.1}{0.9}, from=1-7, to=1-5]
	\arrow["{L_e}", shift right=3, curve={height=-24pt}, between={0.1}{0.9}, from=1-7, to=1-5]
\end{tikzcd}\]

where $\text{res}_e \defeq e(\mhyphen)$ is the \textbf{restriction functor}, and $T_e \defeq \incidencealgebra{P}e \otimes_{\incidencealgebra{P'}} \mhyphen$, $L_e \defeq \Hom_{\incidencealgebra{P'}}(e\incidencealgebra{P}, \mhyphen)$ are the \textbf{idempotent embedding functors}.

\begin{lemma}[{c.f.~\cite[Ch.I, Theorem 6.8]{sim:vol1}}]\label{lemma:ffsubtraction}
	We have the following:
	\begin{enumerate}
		\item $T_e$ and $L_e$ are fully faithful;
		\item $\text{res}_e T_e \cong 1_{\Mod{\incidencealgebra{P'}}} \cong \text{res}_e L_e$ (resp.~$\text{res}_e T_e \cong 1_{\Modf{\incidencealgebra{P'}}} \cong \text{res}_e L_e$);
		\item $T_e \dashv \text{res}_e \dashv L_e$.
	\end{enumerate}
\end{lemma}

Working with representations of the respective bound quivers it is easy to check that, for $M\in \Rep{\HasseQuiver{P'}, I_{P'}} \cong \Mod{\incidencealgebra{P'}}$, we have 
\[\begin{aligned}
	L_e(M)_i &\cong M_i  \quad \forall i=1,\dots, n-1 \quad \text{and}\quad L_e(M)_n \cong \lim_{n<_P j} M_j 
	\\
	T_e(M)_i &\cong M_i  \quad \forall i=1,\dots, n-1 \quad \text{and}\quad T_e(M)_n \cong \colim\limits_{j<_P n } M_j
\end{aligned}
\]
For example, if $M\in \Rep{\HasseQuiver{P'}, I_{P'}}$ and $\widetilde{M} \in \Mod{\incidencealgebra{P'}}$ is the corresponding module we have
\[T_e(M)_i \cong e_i T_e(\widetilde{M}) = e_i \incidencealgebra{P}e \otimes_{\incidencealgebra{P'}} \widetilde{M} \cong e_i \incidencealgebra{P'} \otimes_{\incidencealgebra{P'}} \widetilde{M} \cong e_i \widetilde{M} \cong M_i \qquad \forall i=1,\dots, n-1\]
while $T_e(M)_n \cong e_n T_e(\widetilde{M}) = e_n \incidencealgebra{P}e \otimes_{\incidencealgebra{P'}} \widetilde{M}$. 

We have canonical maps $T_e(M)_j = e_j \incidencealgebra{P}e \otimes_{\incidencealgebra{P'}} \widetilde{M} \xrightarrow{p_{j\to n}\cdot} e_n \incidencealgebra{P}e \otimes_{\incidencealgebra{P'}} \widetilde{M} = T_e(M)_n$ for all $j<_P n$, given by left multiplication with the path from $j$ to $n$ (we have such a path since $j<_P n$ and it is well-defined since all such paths are equal thanks to the relations in the incidence algebra).

Given $e_n p e\otimes m \in T_e(M)_n$ with $p$ a path $\ell \rightsquigarrow n$, we have that $e_n p e\otimes m = p_{\ell\to n}\cdot e_\ell 1_{\incidencealgebra{P}} e \otimes m$ (again, due to the relations in the incidence algebra). With this information, it is easy to check the universal property of the colimit:

\noindent \begin{minipage}{0.55\textwidth}
	given the diagram on the right, we can define (with the previous notation)
	\[g(e_n p e\otimes m) \defeq h_\ell(e_\ell 1_{\incidencealgebra{P}} e\otimes m) \]
	Extending it by linearity, we get the unique map $g\colon T_e(M)_n \to C$ such that $g(p_{j\to n}\cdot (e_j a e)\otimes m) = h_j(e_j a e\otimes m)$ for all $j<_P n$ (thanks to its definition).	
\end{minipage}
\hfill
\begin{minipage}{0.4\textwidth}
	% https://q.uiver.app/#q=WzAsNCxbMCwwLCJUX2UoTSlfaiJdLFsyLDAsIlRfZShNKV9rIl0sWzEsMSwiVF9lKE0pX24iXSxbMSwyLCJDIl0sWzAsMSwicF97alxcdG8ga30gXFxjZG90IiwwLHsic3R5bGUiOnsiYm9keSI6eyJuYW1lIjoiZG90dGVkIn19fV0sWzAsMiwicF97alxcdG8gbn1cXGNkb3QiLDJdLFsxLDIsInBfe2tcXHRvIG59XFxjZG90Il0sWzIsMywiXFxleGlzdHMhIGciLDAseyJzdHlsZSI6eyJib2R5Ijp7Im5hbWUiOiJkb3R0ZWQifX19XSxbMCwzLCJoX2oiLDIseyJjdXJ2ZSI6M31dLFsxLDMsImhfayIsMCx7ImN1cnZlIjotM31dXQ==
	\[\begin{tikzcd}[ampersand replacement=\&,cramped]
		{T_e(M)_j} \&\& {T_e(M)_k} \\
		\& {T_e(M)_n} \\
		\& C
		\arrow["{p_{j\to k} \cdot}", dotted, from=1-1, to=1-3]
		\arrow["{p_{j\to n}\cdot}"', from=1-1, to=2-2]
		\arrow["{h_j}"', curve={height=18pt}, from=1-1, to=3-2]
		\arrow["{p_{k\to n}\cdot}", from=1-3, to=2-2]
		\arrow["{h_k}", curve={height=-18pt}, from=1-3, to=3-2]
		\arrow["{\exists! g}", dotted, from=2-2, to=3-2]
	\end{tikzcd}\]
\end{minipage}
\smallskip


The proof for the functor $L_e$ has a similar procedure.

\begin{definition}
	With the previous setting, we define the \textbf{extension functors} $E^{\Hom}_s$, $E^\otimes_s$ of the subtraction given by $s\colon P'\to P$ as the functors $E^{\Hom}_s\defeq L_e$, $E^\otimes_s\defeq T_e$.
\end{definition}

In the following, we will simply denote them as $E_s$ when we are not be interested in which particular functor we are applying but only on their common properties. Moreover $E_{s \downarrow}$ will denote the restriction to finite-dimensional modules.

